\section{The explicit diagonal and incompleteness proofs}
\label{sec:annex-p05-inc}

\begin{annexitem}{annex:p5.inc.diagonal-coding}{thm:p05-inc-diagonal-coding}{Code construction before self-reference}
\begin{proof}
Use the actual tree constructors in \cref{sec:arithmetic-coding-04}.
The numeral-code function is \(\nu(0)=N_1()\),
\(\nu(n+1)=N_2(\nu(n))\). Induction in the metatheory gives
\(\nu(n)=\ulcorner\bar n\urcorner\). The source's proved
course-of-values traversal recognizes formulas and computes their free
variable lists. Its proved numeral-substitution traversal stops below
an \(x_0\) binder and otherwise commutes with all constructors. Define
\(d(n)\) by that traversal with replacement \(\nu(n)\), when
\(n\) is a formula code with free variables contained in \(\{x_0\}\),
and by the fixed code \(\ulcorner0=0\urcorner\) otherwise. All tests,
the traversal, the numeral recursion and this finite case definition are
primitive recursive by \cref{thm:arithmetization-syntax} and its explicit
closure constructions. Negation is the single constructor \(N_7(n)\)
on formula codes, with \(N_7(\ulcorner0=0\urcorner)\) as its default.
It too is primitive recursive. No representability conclusion follows
from this numerical statement alone.

For a formula \(\beta\) as in the statement, recognition succeeds and
the traversal is precisely the declared capture-avoiding substitution.
The replacement is closed, so all formerly free \(x_0\) occurrences
become closed numeral terms and every other occurrence remains bound.
Thus the result \(G\) is a sentence and its code is exactly \(d(e)\).
Before coding, expand all displayed metalinguistic abbreviations by the
fixed fresh-variable convention. Marked bounds are retained in the
complexity presentation, but their expansion is used for the governing
formula tree. Consequently no extra function symbol, coding sort or
implicit syntactic-equivalence quotient enters this identity.
\end{proof}
\end{annexitem}

\begin{annexitem}{annex:p5.inc.diagonal}{thm:p05-inc-diagonal}{The two directions of the ordinary fixed point}
\begin{proof}
Choose fresh output variable \(y\) and rename bound variables in
\(D\) and \(\theta\) so their use does not capture anything. Here
\(\theta(y)\) means actual capture-avoiding replacement of its only
possible free variable by \(y\). Put
\(\beta(x_0)=\exists y(D(x_0,y)\land\theta(y))\), code its full
expansion as \(e\), and set \(G=\beta(\bar e)\), \(g=d(e)\).
The preceding coding proof gives \(g=\ulcorner G\urcorner\).

The interface at the fixed standard \(e\) proves
\(D(\bar e,\bar g)\), by instantiating its biconditional at
\(\bar g\) and using equality reflexivity. It also proves
\(\forall y(D(\bar e,y)\to y=\bar g)\).
From \(G\), eliminate its existential, apply the latter implication,
and replace equals in \(\theta\), obtaining \(\theta(\bar g)\).
The existential eigenvariable is fresh for the conclusion and all
remaining open assumptions. Conversely, from \(\theta(\bar g)\)
conjoin the proved \(D(\bar e,\bar g)\) and introduce the existential
with witness \(\bar g\). Discharge the displayed hypotheses in each
direction and conjoin the implications. These are first-order deductions
from \(\mathsf Q\); the proof has used no graph claim at nonstandard
inputs and no global normalization of arithmetic formulas.
\end{proof}
\end{annexitem}

\begin{annexitem}{annex:p5.inc.pi1-fixedpoint}{thm:p05-inc-pi1-fixedpoint}{A universal bounded-certificate wrapper}
\begin{proof}
Choose \(y,\vec s\) and all bound variables in the two input formulas
fresh and pairwise distinct where required. In particular \(x_0\)
remains the only possible free variable of \(\beta\). Such a choice
is an effective finite variable-list computation; take successive least
indices absent from the finite union and perform the book's actual
capture-avoiding renamings. Construct the displayed \(\beta\) as an
actual finite arithmetic formula, expand its abbreviations, compute its
code \(e\), and substitute \(\bar e\). By the code proof, the
resulting sentence \(G\) has code \(g=d(e)\).

First check complexity directly in the active hierarchy, not by a
semantic replacement. The fixed negation-normal-form conversion sends
the matrix implication to
\(\neg\delta_d(x_0,y,\vec s)\lor\theta(y)\), with the two
subformulas converted recursively. The converted negative bounded
formula is bounded: De Morgan and marked-bound dualities do not add an
unbounded quantifier. The converted \(\theta\) is \(\Pi^0_1\)
by hypothesis. Closure under disjunction and unbounded universals
therefore puts this exact \(\beta\) in \(\Pi^0_1\). Substitution
of a closed numeral preserves every marked bound's variable condition
and cannot create an unbounded quantifier. Hence \(G\) is
\(\Pi^0_1\). This argument also covers a bounded \(\theta\),
vacuous quantifiers, and a zero-length certificate tuple.

For the biconditional, work in \(\mathsf Q\) using the single graph
instance at \(e\). It gives
\[
 \exists\vec s\,\delta_d(\bar e,\bar g,\vec s),\qquad
 \forall y\bigl((\exists\vec s\,\delta_d(\bar e,y,\vec s))
                                      \to y=\bar g\bigr).
\]
From \(\theta(\bar g)\), fix arbitrary \(y,\vec s\), assume
\(\delta_d(\bar e,y,\vec s)\), introduce its existential tuple,
derive \(y=\bar g\), and use equality replacement to derive
\(\theta(y)\). Discharge the certificate hypothesis and introduce
the universals; their variables are absent from the closed assumptions
\(\mathsf Q\) and \(\theta(\bar g)\). This gives \(G\).
Conversely, from \(G\) instantiate \(y=\bar g\) and eliminate
the proved existential certificate tuple. Its fresh eigenvariables
provide the matrix premise, so modus ponens gives
\(\theta(\bar g)\). The conclusion is independent of the certificate
variables, permitting every existential elimination. Discharge \(G\).
Together these prove the required biconditional.

For the prefix clause, choose \(\vec u\) fresh for \(\delta_d\)
and use the displayed universal prefix at the time \(\beta\) is
formed. Its implication has bounded matrix. To meet the strict
NNF-matrix definition of prefix form, compute its fixed bounded NNF
before coding \(\beta\). This is an explicit finite syntax operation:
on atoms retain the chosen sign, at negation reverse that sign, at
conjunction/disjunction choose the same connective for a positive sign
and its De Morgan dual for a negative sign, replace implication by
the corresponding disjunction/conjunction of its signed children,
and at a marked bounded quantifier retain or dualize that quantifier
according to the sign, recursively converting its matrix. All variable
names and arithmetic bounds remain in place. Every recursive call
visits a proper subformula, so the construction terminates; an
accumulated child-code table computes its two signed versions by
the source's primitive-recursive course-of-values method.
Classical Boolean identities and quantifier dualities prove in
\(\mathsf Q\), by a structural proof-construction induction, that
the resulting bounded matrix is equivalent to the original one.
For marked quantifiers these are the dualities of their actual guarded
expansions; no quantifier is moved through another quantifier. The
result is already prefix-form \(\Pi^0_1\), and \(e\) is the code
of this resulting formula. The forward and
backward deductions just given now derive every
\(\eta(\bar g,\vec u)\) and introduce or instantiate
\(\forall\vec u\) explicitly. Equivalently, the only quantifier
movement involved is the pure-logic equivalence
\(A\to\forall u B\leftrightarrow\forall u(A\to B)\), when
\(u\) is absent from \(A\); its two directions respectively use
universal elimination followed by introduction, and those same rules
in reverse order under the temporary assumption \(A\).
There is no move of an unbounded witness through a bounded universal.
The \(\mathcal N\)-equivalence of general prefix normal forms in
\cref{thm:prefix-normal-form} is neither needed nor promoted to a
\(\mathsf Q\)-provable equivalence.
\end{proof}
\end{annexitem}

\begin{annexitem}{annex:p5.inc.proof-facts}{thm:p05-inc-proof-facts}{Numeral evidence and the two consistency implications}
\begin{proof}
If \(T\derives\sigma\), actual certified-proof adequacy in
\cref{thm:arithmetization-syntax} supplies a standard proof code
\(p\) for \(a=\ulcorner\sigma\urcorner\). Positive numeralwise
representation gives a \(\mathsf Q\)-proof of
\(P_T(\bar p,\bar a)\); it is also a \(T\)-proof since
\(T\supseteq\mathsf Q\). Existential introduction with \(\bar p\)
proves \(\operatorname{Pr}_T(\bar a)\). This is the external first
derivability condition, not a universal internal proof transformer.
In \(\mathcal N\), the declared defining formula and standard-domain
existential semantics give
\(\operatorname{Pr}_T(\bar a)\) iff an actual numerical proof code
exists. Adequacy gives exactly \(T\derives\sigma\).

All seven \(\mathsf Q\) axioms hold in \(\mathcal N\): successor
is injective and nonzero, every nonzero natural number has a predecessor,
and standard addition and multiplication satisfy the four defining
equations. The actual natural-deduction soundness theorem
\cref{thm:fol-soundness}, for the active calculus rather than an
off-route rule convention, therefore transfers each \(\mathsf Q\)
derivation to truth in \(\mathcal N\). In
particular \(\mathsf Q\) proves \(0\neq S0\), and \(0=S0\)
is a false bounded, hence \(\Sigma^0_1\), sentence. An inconsistent
\(T\) proves it by explosion, contradicting \(\Sigma^0_1\)-soundness.
For omega-consistency, an inconsistent \(T\) proves
\(\exists x(x=x)\) and, for each standard \(n\), proves
\(\neg(\bar n=\bar n)\). This is precisely a prohibited
omega-inconsistency witness. Thus each stronger hypothesis implies
consistency. Neither direction compares the two stronger hypotheses.
\end{proof}
\end{annexitem}

\begin{annexitem}{annex:p5.inc.godel-chain}{thm:p05-inc-godel-chain}{Consistency, standard truth and the two nonrefutation branches}
\begin{proof}
The proof formula \(P_T\) is \(\Sigma^0_1\); prepending
\(\exists p\) keeps that class. By the fixed NNF conversion,
its negation is \(\Pi^0_1\): De Morgan interchanges the two
connectives, bounded duality stays bounded, and each unbounded
existential becomes universal. This is an induction on the generating
clauses of the syntactic class, not an arithmetic collection theorem.
The universal certificate construction consequently produces the
stated \(G_T\), with its biconditional already in \(\mathsf Q\).

Assume consistency. If \(T\derives G_T\), proof evidence gives
\(T\derives\operatorname{Pr}_T(\bar g)\), while the biconditional
gives \(T\derives\neg\operatorname{Pr}_T(\bar g)\). This violates
consistency, so \(T\notderives G_T\). Consequently there is no
standard proof code for \(g\), and proof adequacy in \(\mathcal N\)
gives \(\mathcal N\entails\neg\operatorname{Pr}_T(\bar g)\).
The biconditional is a \(\mathsf Q\) theorem and hence true in
\(\mathcal N\). Therefore \(\mathcal N\entails G_T\).
There is no soundness assumption about all axioms of \(T\) here.

Assume \(\Sigma^0_1\)-soundness. It implies consistency by the
preceding lemma, so the preceding branch shows that \(G_T\) is true.
Its negation is syntactically \(\Sigma^0_1\), by the same dual
generating-clause argument, and false in \(\mathcal N\).
Thus \(T\) cannot prove \(\neg G_T\).

Assume omega-consistency. Again obtain consistency first. If
\(T\derives\neg G_T\), classical logic applied to the fixed-point
biconditional gives \(T\derives\operatorname{Pr}_T(\bar g)\),
that is, \(T\derives\exists p\,P_T(p,\bar g)\).
By the first branch \(T\notderives G_T\), so each standard number
\(n\) fails the numerical proof check for \(g\). The negative
numeralwise clause provides a \(\mathsf Q\)-proof, hence a
\(T\)-proof, of \(\neg P_T(\bar n,\bar g)\) for every \(n\).
The single formula \(P_T(p,\bar g)\), with at most free variable
\(p\), is then an omega-inconsistency witness. Contradiction.
This argument does not infer an internal universal rejection from all
the separate standard numeral rejections.

For the literal prefix selection
\(P_T(p,a)=\exists\vec w\,\kappa_T(p,a,\vec w)\), pure first-order
quantifier duality proves
\(\neg\operatorname{Pr}_T(a)\leftrightarrow
\forall p\,\forall\vec w\,\neg\kappa_T(p,a,\vec w)\)
in \(\mathsf Q\). There is no bounded universal across which the
witnesses must be collected. Feed this literal prefix formula into
the prefix construction, and compute \(e\) as the code of that
particular resulting \(\beta\), not of its unprefixed alternative.
The code proof still gives the exact self-substitution identity.
Before coding this prefix \(\beta\), use the same explicit signed
bounded NNF construction on the matrix
\(\neg\delta_d\lor\neg\kappa_T\). Thus the actual coded prefix
has a bounded NNF matrix, while the statement's implication is its
declared logically equivalent display. This supplies the precise
prefix form claimed. The three proof branches are unchanged because
their biconditional is with the same selected \(\operatorname{Pr}_T\).
\end{proof}
\end{annexitem}

\begin{annexitem}{annex:p5.inc.rosser-chain}{thm:p05-inc-rosser-chain}{The fixed point and both bounded cases}
\begin{proof}
Write \(H(a,b)\) for the actual formula
\(\forall y(P_T(y,a)\to\exists z\leq y\,P_T(z,b))\).
The letters \(a,b\) are arithmetic variables; the expression
\(\operatorname{neg}(a)\) is not used as an arithmetic term.
Choose fresh variables and form
\[
 \beta(x_0)=\exists a\,\exists b\,
                   (D(x_0,a)\land N(a,b)\land H(a,b)).
\]
Compute \(e=\ulcorner\beta\urcorner\), set \(R_T=\beta(\bar e)\),
\(r=d(e)=\ulcorner R_T\urcorner\), and
\(r^-=\operatorname{neg}(r)=\ulcorner\neg R_T\urcorner\).
The graph instances at the two standard inputs \(e\) and \(r\)
prove \(D(\bar e,\bar r)\), \(N(\bar r,\overline{r^-})\),
and uniqueness of each output. Eliminating the two existentials
and replacing equals proves \(R_T\to H(\bar r,\overline{r^-})\).
Conversely conjoin those two proved graph facts with
\(H(\bar r,\overline{r^-})\) and introduce the existential witnesses
\(\bar r,\overline{r^-}\). This proves the exact displayed
biconditional in \(\mathsf Q\), without a partial negation function
abbreviation or an unproved internal graph composition.

First suppose \(T\derives R_T\), and choose an actual standard
proof code \(p\). The positive numeralwise clause gives
\(T\derives P_T(\bar p,\bar r)\). The fixed-point implication,
universal elimination at \(\bar p\), and modus ponens give
\[
 T\derives\exists z\leq\bar p\,P_T(z,\overline{r^-}).
\]
Consistency implies there is no numerical proof of \(\neg R_T\).
For every standard \(i=0,\ldots,p\), the negative numeralwise
clause therefore provides an actual \(\mathsf Q\)-derivation of
\(\neg P_T(\bar i,\overline{r^-})\).
The existing \(p\)-instance of \cref{lem:q-bounded-cases} is a
\(\mathsf Q\)-proof of
\(z\leq\bar p\to\bigvee_{i=0}^p z=\bar i\).
Under \(z\leq\bar p\), each equality disjunct substitutes into
the relevant negative proof, giving \(\neg P_T(z,\overline{r^-})\).
Finite disjunction elimination, implication discharge and universal
introduction yield
\[
 \mathsf Q\derives\forall z
       (z\leq\bar p\to\neg P_T(z,\overline{r^-})).
\]
Existential elimination refutes the displayed bounded existential
already derived in \(T\), contradicting consistency. A least proof
code was not required. The finite conjunction of separate numeral
rejections is justified by actual disjunction elimination, not by
an assertion that every standard number is one of these numerals.

Now suppose \(T\derives\neg R_T\), with actual standard proof
code \(q\). Positive representation gives
\(\mathsf Q\derives P_T(\bar q,\overline{r^-})\).
Consistency implies there is no numerical proof of \(R_T\), so
negative representation gives, for each \(i=0,\ldots,q\),
\(\mathsf Q\derives\neg P_T(\bar i,\bar r)\).
The \(q\)-instance of bounded cases and equality replacement,
as in the preceding branch, now give
\[
 \mathsf Q\derives\forall y
                   (y\leq\bar q\to\neg P_T(y,\bar r)).
\]
Fix arbitrary \(y\) and temporarily assume \(P_T(y,\bar r)\).
The existing \(q\)-instance of \cref{lem:q-dichotomy} gives
\(y\leq\bar q\lor\bar q\leq y\). In the first disjunct the
last negative implication contradicts the temporary assumption;
explosion proves \(\exists z\leq y\,P_T(z,\overline{r^-})\).
In the second disjunct, conjoin \(\bar q\leq y\) with the
proved positive \(P_T(\bar q,\overline{r^-})\), and introduce
that bounded existential with witness \(\bar q\).
Disjunction elimination gives the same conclusion in both branches.
Discharge \(P_T(y,\bar r)\) and introduce \(\forall y\): all
other open assumptions are closed \(\mathsf Q\) axioms.
Thus \(\mathsf Q\), and hence \(T\), proves
\(H(\bar r,\overline{r^-})\); the reverse fixed-point implication
gives \(T\derives R_T\), contradicting its supposed proof of
\(\neg R_T\). Both unprovability claims are established using
consistency alone, but for this different sentence.

The two comparison lemmas are individually derived by metainduction
on the fixed standard bound, using only successor, predecessor cases,
and the two right-recursion addition axioms. No object-language
induction, commutativity of addition, global order linearity or
internal minimal-proof principle was used here. The result remains
valid for the stated \(\Sigma^0_1\) numeralwise proof formula;
no \(\Delta^0_0\) classification of that formula was assumed.
\end{proof}
\end{annexitem}

\begin{annexitem}{annex:p5.inc.finite-extension}{thm:p05-inc-finite-extension}{The exact finite deduction and consistency test}
\begin{proof}
Take \(\bot\) as the contradiction formula of the actual calculus.
If \(T\cup D\) is inconsistent, its two contradictory derivations
combine by \(\neg E\) to a finite derivation of \(\bot\).
Replace every open leaf labelled by a member of the finite list \(D\)
by a derivation of that member from a single temporary closed
assumption \(C_D\), using the appropriate finite conjunction
eliminations. Other open leaves still belong to \(T\).
Repeated occurrences use copies, not shared derivation rows.
The replacement assumption is closed, so it creates no new free
variable violating \(\forall I\). Closedness alone would not meet
the source's stronger \(\exists E\) condition: a closed conjunction
can still contain the eigenvariable as a bound variable. Before leaf
replacement, rename each existential-elimination eigenvariable to an
index absent from the all-variable lists of the entire finite proof
and \(C_D\). In its minor derivation, replace its free occurrences,
including those in terms and rule parameters, by that fresh variable.
An induction on that minor derivation verifies this operation. In
propositional rules the replacement distributes over the formula
constructors. In the two instantiation rules and equality replacement,
term/formula substitution commutes with the free renaming up to alpha
conversion, which the actual alpha rule supplies. At a binder with
the old variable name, free replacement stops. A universal-introduction
step on that variable is rebuilt with the fresh variable instead:
its renamed premise yields the renamed bound formula, followed by
an alpha step to the required conclusion. On another variable its
side condition is preserved. A nested existential elimination is
treated with its own fresh eigenvariable; the chosen fresh names
cannot clash with any binder, rule parameter or other eigenvariable.
The designated witness assumption becomes exactly the freshly
instantiated matrix. Remaining open assumptions and the conclusion
had no occurrence of the old eigenvariable by the original side
condition and therefore remain unchanged. Bound-variable alpha
renamings, when required by a substitution parameter, use the actual
alpha rule. Every \(\exists E\) application now satisfies its
all-occurrence condition also for \(C_D\). Choose fresh discharge
labels for the replacement occurrences, so the replacement does not
accidentally discharge another premise branch's occurrence.

The resulting derivation of \(\bot\) from \(T\) and \(C_D\)
gives \(T\derives\neg C_D\) by \(\neg I\). If no replaced
leaf was present, the book's nonvacuous discharge still applies after
adjoining a use of the temporary assumption: from \(\bot\) infer
\(\neg C_D\) by explosion; conjoin or use it with the assumed
\(C_D\) to infer \(\bot\) by \(\neg E\); then discharge that
actual occurrence of \(C_D\). For the empty list, use the actual
proof of \(0=0\) in place of any conjunction eliminations. Thus the
argument also covers empty \(D\), repetitions in the list, and a
contradiction already provable in \(T\).

Conversely, if \(T\derives\neg C_D\), assumptions from the list
\(D\) give \(C_D\) by successive \(\land I\), or \(=I\)
gives \(0=0\) when the list is empty. With the given negative
derivation, \(\neg E\) gives \(\bot\). Explosion yields a
sentence and its negation, so \(T\cup D\) is inconsistent.
This proves the biconditional for the actual calculus, rather than
invoking a deduction theorem from an unspecified Hilbert system.

Finally put \(S=T\cup D\), assumed consistent. If \(S+\sigma\)
is inconsistent, the same argument with singleton list \(\sigma\)
gives \(S\derives\neg\sigma\). If also \(S+\neg\sigma\)
is inconsistent, it gives \(S\derives\neg\neg\sigma\), hence
\(S\derives\sigma\) by the actual classical \(\mathsf{RAA}\)
rule: assume \(\neg\sigma\), derive contradiction with the double
negation, and discharge that occurrence. This contradicts consistency
of \(S\). Therefore at least one extension is consistent.
\end{proof}
\end{annexitem}

\begin{annexitem}{annex:p5.inc.theoremhood-chain}{thm:p05-inc-theoremhood-chain}{A computed complete extension and its Rosser contradiction}
\begin{proof}
Certified numerical proof checking is primitive recursive, with
sentence recognition primitive recursive, by
\cref{thm:arithmetization-syntax}. Enumerate all pairs \((p,a)\)
by finite increasing-square stages and output \(a\) exactly when
it is a sentence code passing that proof check. Every pair appears
eventually, each stage ends, and proof adequacy shows that the outputs
are precisely the theorem codes. Thus sentence theoremhood is c.e.

Suppose now a total decider for this set exists. Enumerate all
sentences \(\sigma_0,\sigma_1,\ldots\) by increasing valid sentence
code. Recognition is primitive recursive and there are infinitely
many sentences, for example numeral equalities; finding each next
code terminates. Maintain a finite list \(D_n\), starting empty.
At stage \(n\), form the actual sentence
\(\neg C_{D_n\mathbin{+}[\sigma_n]}\) using the constructor
codes and query the supposed \(T\)-theoremhood decider on its code.
If the answer is no, append \(\sigma_n\) to \(D_n\); if yes,
append \(\neg\sigma_n\). This is a total computable recursion
on finite lists, not a search through consistency proofs.

Prove in the metatheory that each \(S_n=T\cup D_n\) is consistent.
The base is the assumed consistency of \(T\). In the no-answer
case, the finite extension biconditional shows exactly that
\(S_n+\sigma_n\) is consistent, so this is \(S_{n+1}\).
In the yes-answer case it shows exactly that
\(S_n+\sigma_n\) is inconsistent. Since \(S_n\) is consistent,
the last clause of the finite extension lemma proves that
\(S_n+\neg\sigma_n\), the selected \(S_{n+1}\), is consistent.
This induction is an external proof about the computed construction;
it does not insert an induction axiom into any \(S_n\) or into
\(\mathsf Q\).

Let \(T^*=T\cup\bigcup_nD_n\). Dovetail the c.e. enumeration of
the original sentence axiom set \(T\) with the computed finite
additions. This enumerates a c.e. axiom set \(T^*\supseteq\mathsf Q\).
The computability of the list sequence is not incorrectly identified
with decidability of membership in its union. If \(T^*\) were
inconsistent, a finite derivation of contradiction would use finitely
many of the added axiom occurrences. Their finitely many stages have
a maximum \(N\), so the same derivation is from \(T\cup D_N\),
contradicting the proved consistency of \(S_N\). Hence \(T^*\)
is consistent. For every sentence \(\sigma\), its enumeration
stage adjoins \(\sigma\) or \(\neg\sigma\). That axiom itself
is a derivation from \(T^*\). Thus \(T^*\) is deductively complete.

If the phrase ``computable complete extension'' denotes a decidable
theory, its deductive closure has that additional property as well:
on any sentence input \(\tau\), enumerate \(T^*\)-proofs in
parallel until a proof of \(\tau\) or a proof of \(\neg\tau\)
appears. Completeness ensures termination; consistency ensures that
the two answers cannot both occur. The affirmative answer is precisely
membership in its deductive closure, by proof adequacy. This algorithm
requires only the c.e. axiom enumeration, with the usual finite axiom
certificates, not decidable membership in the axioms.

Apply \cref{thm:p05-inc-rosser-chain} to the c.e., consistent
\(\mathsf Q\)-extension \(T^*\). Its required certificate
interface is instantiated for this axiom enumeration, as quantified
in the present theorem's hypotheses. It supplies a sentence \(R_{T^*}\)
whose two sides are both underivable there. Deductive completeness
gives one side, a contradiction. Therefore the original theoremhood
decider cannot exist. A theorem set for arbitrary formulas would also
decide its subset of sentence codes by restriction, so that broader
version is undecidable too. Merely observing that a decider might
answer no to both sides of a Rosser pair was not the argument: the
decider was used to compute the forbidden complete extension.
\end{proof}
\end{annexitem}
